% Fragmento da Metodologia para o template SBC (Overleaf).
% Issue #30 — colar após a Introdução. Atualize a tabela do funil com funnel.csv.

\section{Metodologia}

\subsection{Fonte de dados e janela de observação}

A unidade de análise é o repositório público no GitHub com GitHub Actions.
Os dados vêm da API REST pública, com cache local, \textit{rate limit} e
\textit{backoff} exponencial; o token permanece apenas em variável de ambiente.

A janela de observação é de 12 meses (UTC), configurada em \texttt{config.yaml}.
Nesta sprint: 2025-10-01 a 2026-09-30. Contagens e métricas temporais consideram
eventos com \textit{timestamp} na janela, salvo âncoras permitidas pelo enunciado
(ex.: release anterior fora da janela no \textit{compare}).

\subsection{Funil de seleção}

A busca é fatiada por faixas de estrelas (teto de 1.000 resultados por consulta).
Critérios de inclusão: repositório não arquivado; $\geq 5$ releases publicadas
(\texttt{draft=false}, \texttt{prerelease=false}) na janela; $\geq 50$
\textit{workflow runs} \texttt{push} no \textit{default branch} na janela.
Meta Lab03S02: $\geq 300$ repositórios. Com estratificação por faixa
(\texttt{sample.stratify}), o tamanho-alvo é dividido em cotas iguais ---
escolha que super-representa faixas altas e deve constar no artigo.
A tabela quantitativa do funil é gerada em \texttt{funnel.csv}.

\subsection{Definições operacionais e variantes}

\begin{itemize}
  \item \textbf{Frequência de deploy:} releases na janela $/$ semanas da janela;
    variantes com pré-releases e \textit{tags} (RQ\,07).
  \item \textbf{Lead time:} (a) mediana por release do commit mais antigo do
    \textit{compare}; (b) mediana sobre todos os commits.
  \item \textbf{CFR:} (a) falhas de CI; (b) releases seguidas em $\leq 7$ dias
    por release corretiva (patch SemVer ou \texttt{revert}/\texttt{hotfix}/\texttt{fix});
    releases nos 7 dias finais são censuradas.
  \item \textbf{Recuperação:} mediana dos episódios concluídos no mesmo
    \textit{workflow}; episódios sem sucesso são censurados; repos sem episódio
    ficam com valor ausente.
\end{itemize}

Tempos em horas; proporções em $[0,1]$. A classificação DORA usa os cortes da
disciplina; a nota geral é a mediana das quatro notas com arredondamento para
baixo. A referência \textbf{C1} combina release, lead time (a), CFR (a) e
recuperação.

\subsection{Protocolo de validação manual}

Amostra-ouro: 60 repositórios e 5 releases cada, sementes fixas
(\texttt{random\_state\_repos=42}; releases $42+$índice). Três avaliadores
rotulam independentemente tipo de projeto, entrega real e se cada release é
corretiva, \textbf{sem} ver a heurística. Concordância: kappa de Fleiss
interpretado por Landis \& Koch. Consenso: maioria simples; empates 1-1-1 em
reunião. Sobre o consenso, precisão/\textit{recall}/F1 da heurística (e
possível refinamento).

\subsection{Ameaças à validade (prévia)}

Construto (proxy $\neq$ produção), interna (censuras), externa (GitHub Actions
e estratificação) e de conclusão (múltiplas variantes e correção de Holm).
